\documentclass[12pt, letterpaper]{article}

% ── Packages ────────────────────────────────────────────────────────────────
\usepackage[margin=1in]{geometry}
\usepackage{parskip}
\usepackage{enumitem}
\usepackage[hidelinks]{hyperref}
\usepackage{titlesec}
\usepackage{booktabs}
\usepackage{array}
\usepackage{tabularx}
\usepackage{graphicx}
\usepackage{fancyhdr}
\usepackage{xcolor}
\usepackage{setspace}
\usepackage{microtype}
\usepackage{caption}
\usepackage{mdframed}
\usepackage{listings}
\usepackage{fontenc}
\usepackage{inputenc}

% ── Typography ───────────────────────────────────────────────────────────────
\singlespacing
\setlength{\parindent}{0pt}
\setlength{\parskip}{8pt}

\titleformat{\section}{\large\bfseries}{}{0em}{}[\vspace{-4pt}\rule{\linewidth}{0.6pt}\vspace{2pt}]
\titleformat{\subsection}{\normalsize\bfseries}{}{0em}{}
\titleformat{\subsubsection}{\normalsize\itshape}{}{0em}{}

% ── Header / Footer ──────────────────────────────────────────────────────────
\pagestyle{fancy}
\fancyhf{}
\rhead{\small International Student AI Influencer}
\lhead{\small Gen AI \& Social Media}
\cfoot{\thepage}
\renewcommand{\headrulewidth}{0.4pt}

% ── Listings style ───────────────────────────────────────────────────────────
\lstset{
  basicstyle=\ttfamily\footnotesize,
  breaklines=true,
  frame=single,
  backgroundcolor=\color{gray!10}
}

% ─────────────────────────────────────────────────────────────────────────────
\begin{document}

% ── Title Page ───────────────────────────────────────────────────────────────
\begin{titlepage}
  \centering
  \vspace*{2cm}
  {\LARGE\bfseries International Student AI Influencer\par}
  \vspace{0.5cm}
  {\large An AI-Powered Multi-Agent Pipeline for Civic Content Generation\par}
  \vspace{2cm}
  {\large\itshape Gen AI \& Social Media\par}
  \vspace{1.5cm}

  \begin{tabular}{c}
    \large Joaquin Sardon \\[4pt]
    \large Kulbir Kaur \\[4pt]
    \large Melissa Lee \\[4pt]
    \large Shivani Kabra \\[4pt]
    \large Gerald Velasquez \\
  \end{tabular}

  \vspace{2cm}
  {\large May 2026\par}
  \vfill
  {\small NYC Build With AI Hackathon 2026}
\end{titlepage}

\newpage
\thispagestyle{empty}
\mbox{}
\newpage

% ── Executive Summary ────────────────────────────────────────────────────────
\section*{Executive Summary}
\addcontentsline{toc}{section}{Executive Summary}

The International Student AI Influencer is a multi-agent generative AI pipeline that
transforms fragmented civic and financial information into platform-ready social media
content targeting international students in the United States. Built during the NYC
Build With AI Hackathon 2026, the system addresses a concrete and underserved gap: no
single accessible channel consolidates the housing, tax, immigration, banking, employment,
and safety information that international students urgently need upon arrival. The
pipeline automates the full content lifecycle---source retrieval from live web results,
NYC Open Data, and curated RSS feeds; post drafting; AI image generation; five-scene
video production with lip-synced voiceover; and automated quality scoring---and surfaces
it through a React-based web interface backed by a FastAPI server. The main takeaway
is that a compact, well-chained set of specialized AI agents can bridge the gap between
raw civic data and polished, student-appropriate social media content at low marginal
cost, while making explicit the trust and accuracy risks that any such system must
eventually resolve before production deployment.

\textbf{GitHub Repository:}
\url{https://github.com/gerald512-beep/InterStudent}

\newpage

% ─────────────────────────────────────────────────────────────────────────────
\section{Introduction}
% ─────────────────────────────────────────────────────────────────────────────

\subsection{Problem Statement and Motivation}

More than one million international students study at colleges and universities in the
United States each year. Upon arrival---and repeatedly throughout their enrollment---they
must navigate a dense, poorly integrated web of civic systems: opening a bank account
without a Social Security Number, understanding F-1 visa employment restrictions,
filing taxes as a non-resident alien, securing housing without a US credit history,
accessing city emergency services, and interpreting rapidly changing immigration policy.
The information that governs each of these domains lives in separate silos: IRS
publications, university international student offices, state housing authority portals,
local government APIs, and financial institution blogs. No unified, student-facing
channel aggregates and translates this information into a format appropriate to where
students actually consume media: social platforms.

The consequence is predictable. Students rely on informal peer networks, outdated
Reddit threads, or advice from friends who faced the same uncertainty a year or two
earlier. The information they act on is frequently wrong, incomplete, or dangerously
out of date---particularly in domains like tax law and immigration compliance where
errors carry serious penalties.

This project asks: can a generative AI pipeline reliably bridge the gap between raw
civic data sources and student-appropriate social media content, at a quality and
speed that a small team or student organization could sustain?

\subsection{Who the Project Serves}

The primary audience is international students in the United States, specifically those
on F-1 and J-1 visas navigating financial and civic systems for the first time. The
prototype targets New York City and draws on NYC-specific datasets, but the pipeline
architecture is designed to be portable: swapping the data connectors in Agent 1 would
allow the same pipeline to serve international student communities at any US university
or in any city with comparable open data infrastructure.

A secondary audience is the student organizations, international student offices, and
nonprofit groups that could operate the tool as a content channel---producing a regular
cadence of accurate, timely, platform-native posts without requiring a dedicated
communications team.

\subsection{Generative AI Angle}

The project uses generative AI not as a chatbot or search interface, but as a
coordinated production pipeline. Four specialized agents plus an output-rendering
layer each apply a different model or modality at the appropriate stage of content
creation:

\begin{itemize}[leftmargin=1.5em]
  \item \textbf{Semantic retrieval} uses text embeddings (Google's
    \texttt{gemini-embedding-001}) to rank heterogeneous documents by relevance to
    a given civic topic, with topic-specific weight boosting.
  \item \textbf{Content drafting} uses a large language model (Gemini 2.5 Flash) to
    detect the trending angle in retrieved data and generate post text, hashtags,
    and a five-scene video storyboard with an SSML voice script.
  \item \textbf{Image generation} uses Imagen 4 Ultra (with a fallback chain to
    Imagen 4 and Imagen 3) to produce a photorealistic social image aligned to the
    post and the avatar persona.
  \item \textbf{Video production} uses Veo 3.1 Fast to generate per-scene video
    clips with embedded voiceover, assembled by MoviePy into a 9:16 MP4 with
    Chirp3-HD Text-to-Speech audio replacement for precise SSML-controlled pacing.
  \item \textbf{Automated quality control} uses Gemini 2.5 Flash to score the
    output against seven criteria before it enters the publishing queue.
\end{itemize}

Generative AI is the right fit here because the core task---reading unstructured
civic information and producing polished, platform-appropriate media---is exactly what
modern foundation models do well. Rule-based systems or templated content would either
miss the nuance of changing policy or produce content too generic to be useful. The
agentic pattern allows each model to specialize: retrieval, writing, image, video,
and evaluation are distinct cognitive tasks that benefit from being handled separately
and then composed.

\subsection{Limitations of the Status Quo}

Existing solutions fail this audience in different ways. General-purpose AI assistants
(ChatGPT, Gemini, Claude) can answer individual questions but do not monitor live
data sources, produce platform-formatted posts automatically, or maintain a consistent
persona and publishing cadence. Government portals and university international student
office websites are authoritative but static, text-heavy, and not distributed through
the channels where students spend time. Commercial social media management tools
(Buffer, Hootsuite, Sprout Social) provide scheduling and analytics but have no native
ability to generate civic content from raw open data. The gap between raw information
and student-ready media is currently filled by human effort---or not filled at all.

\subsection{Ethics, Risk, and Governance}

The project team explicitly identified four categories of risk during design, and this
report treats them as first-class concerns rather than an afterthought.

\textbf{Accuracy and currency.} Generated posts derive from live data sources, but
the pipeline cannot verify that retrieved information is correct or up to date before
publishing. An incorrect statement about tax filing deadlines or visa employment
restrictions could cause real harm to a student who acts on it.

\textbf{Overgeneralization.} The system uses a single audience persona configuration.
International students represent more than 200 nationalities, dozens of visa categories,
and highly varied prior experience with US institutions. A post calibrated for one
segment may be misleading or irrelevant to another.

\textbf{Implicit advice on sensitive topics.} Housing, immigration, and financial
decisions carry legal and financial consequences. Content that presents itself as
informational can function as advice in practice, particularly when it is polished,
confident, and platform-native in presentation.

\textbf{Content quality theater.} A post can be visually polished, grammatically
correct, and structurally on-brand while being substantively too generic to be
useful---or subtly wrong in ways that are not apparent to a student reading it for
the first time.

These risks motivate the roadmap described in Section 4 and are directly reflected in
the design of Agent 4 (quality control) and the decision to implement a human approval
gate before any content is published.

\newpage

% ─────────────────────────────────────────────────────────────────────────────
\section{Implementation}
% ─────────────────────────────────────────────────────────────────────────────

\subsection{System Overview}

The system comprises five layers that correspond to five functional concerns: data
retrieval, content generation, image rendering, video production, and quality control.
These layers are implemented as discrete Python modules orchestrated by a FastAPI
backend and surfaced to users through a React single-page application. Figure 1 shows
the high-level pipeline.

\begin{mdframed}[backgroundcolor=gray!8, linewidth=0.6pt]
\begin{center}
\small
\textbf{Figure 1. High-Level Agent Pipeline}\\[6pt]
\texttt{
\begin{tabular}{ccccccc}
[User Input] & $\to$ & Agent 1 & $\to$ & Agent 2 & $\to$ & Creative Storyteller \\
             &        & Retrieval &      & Content &       & Image + Final Text \\
             &        &         &      & Generator &     & \\[6pt]
             &        &         &      & $\downarrow$ & & $\downarrow$ \\[2pt]
             &        &         &      & Agent 3 & $\to$ & Agent 4 \\
             &        &         &      & Video   &       & QC + Decision \\
             &        &         &      &         &       & $\downarrow$ \\
             &        &         &      &         &       & Publishing Queue \\
\end{tabular}
}
\end{center}
\end{mdframed}

Agent 5 (Scenario Resolver) operates as a standalone optional tool for deterministic
financial guidance and does not participate in the main content pipeline.

\subsection{Agent 1: Source Retrieval}

Agent 1 is the data foundation of the pipeline. Given a civic topic query, it
aggregates documents from four source categories:

\begin{itemize}[leftmargin=1.5em]
  \item \textbf{Google Search Grounding via Gemini 2.5 Flash}: Live web results
    with source citations. This is the highest-signal input because it reflects
    current policy, news, and guidance rather than a static snapshot.
  \item \textbf{NYC Open Data (Socrata API)}: Three curated datasets---NYC job
    postings (\texttt{kpav-sd4t}), CUNY enrollment figures (\texttt{d44m-4xgv}),
    and NYC living wage benchmarks (\texttt{ic3t-wcy2}).
  \item \textbf{Financial RSS Feeds}: Eight feeds spanning IRS announcements, CFPB
    consumer updates, NASFAA financial aid news, NYT Business, and four student
    loan providers (SoFi, Prodigy Finance, Earnest, MPOWER).
  \item \textbf{Static Tax Knowledge Seed}: Seven curated F-1/J-1 tax scenarios
    hardcoded as a fallback for domains where live retrieval is unreliable.
\end{itemize}

Retrieved documents are chunked at 400 characters with a 50-character overlap and
embedded using \texttt{gemini-embedding-001}. Relevance ranking uses in-memory cosine
similarity with topic-specific weight boosting, producing a ranked \texttt{retrieval\_pack}
passed downstream. No persistent vector database is required, keeping the architecture
stateless for hackathon purposes.

\subsection{Agent 2: Content Generator}

Agent 2 receives the \texttt{retrieval\_pack} and produces all the raw material for
a social media post. It operates in three stages:

\begin{enumerate}[leftmargin=1.5em]
  \item \textbf{Trend Detection}: Identifies the most timely and engaging angle
    from the retrieved results, along with an urgency signal for CTA framing.
  \item \textbf{Post Drafting}: Generates post text calibrated to the selected
    platform (LinkedIn or Instagram), including hashtags, calls-to-action, and a
    image generation prompt, using persona configuration from
    \texttt{config/persona.py}.
  \item \textbf{Video Brief Generation}: Produces a five-scene storyboard with per-scene
    visual prompts, character direction, and an SSML-annotated voice script designed
    for Chirp3-HD Text-to-Speech.
\end{enumerate}

All inference in Agent 2 runs against Gemini 2.5 Flash. The output is a
\texttt{content\_draft} dictionary consumed by both the Creative Storyteller and
Agent 3 in parallel.

\subsection{Creative Storyteller: Image Rendering and Final Text}

The Creative Storyteller module handles image generation and the final multimodal
polish pass. It implements an Imagen fallback chain: Imagen 4.0-ultra $\to$ Imagen
4.0 $\to$ Imagen 3.0, catching API errors at each step. Once an image is produced,
the image bytes and draft post text are passed together to Gemini 2.5 Flash in a
multimodal prompt, producing a final post body that is tonally and visually coherent
with the generated image. The output is a \texttt{final\_output} dictionary containing
the post body, base64-encoded image, platform, hashtags, and cited sources.

\subsection{Agent 3: Video Generator}

Agent 3 is the most computationally intensive component. Given the five-scene storyboard
from Agent 2, it:

\begin{enumerate}[leftmargin=1.5em]
  \item Synthesizes SSML voiceover audio using Google Cloud Text-to-Speech Chirp3-HD
    (\texttt{en-US-Chirp3-HD-Aoede} for female, \texttt{en-US-Chirp3-HD-Orus} for male).
  \item Generates one Veo 3.1 Fast video clip per scene (9:16, 8 seconds), with the
    voiceover embedded in the generation prompt for natural lip sync and body language.
  \item Runs scenes in parallel using a \texttt{ThreadPoolExecutor} with three workers,
    reducing wall time from approximately 10 minutes sequential to approximately 3 minutes.
  \item Falls back to Imagen 4 Ultra static images with a Ken Burns zoom effect when
    Veo generation fails for a scene.
  \item Assembles the clips into a single MP4 using MoviePy, overlaying subtitle text
    and replacing Veo's native audio track with the Chirp3-HD track to preserve
    SSML-controlled pacing (\texttt{<emphasis>}, \texttt{<prosody>}, \texttt{<break>} tags).
\end{enumerate}

The final video is a 9:16 portrait MP4 at 24 FPS, suitable for Instagram Reels and
TikTok. The base64-encoded video and a thumbnail are returned in the
\texttt{video\_output} dictionary.

\subsection{Agent 4: Quality Control}

Agent 4 uses Gemini 2.5 Flash to evaluate the combined output of the Creative
Storyteller and Agent 3 against seven criteria: accuracy, tone, engagement, compliance,
platform fit, video script coherence, and visual-audio alignment. Each criterion is
scored on a 1--10 scale. An aggregate score at or above 7 triggers a \texttt{PUBLISH}
decision; below 7 the agent returns a \texttt{MODIFY} decision with structured feedback
for each failing criterion. A JSON parse failure defaults to \texttt{PUBLISH} at score
7 to prevent the pipeline from blocking indefinitely.

\subsection{Agent 5: Scenario Resolver}

Agent 5 is an optional, standalone module for deterministic financial guidance. It
implements a two-stage extraction pattern: a language model first extracts the
structured parameters of a user query (visa type, income type, filing status), and
a rule engine (\texttt{scenario\_rules.py}, 142 lines) then maps those parameters to
a deterministic guidance response. This keeps the most legally sensitive domain---tax
and financial advice---out of the probabilistic generation path.

\subsection{Backend: FastAPI Server}

The FastAPI server (\texttt{backend/server.py}) wraps the agent pipeline in a REST
interface. Key endpoints are listed below.

\begin{center}
\begin{tabularx}{\textwidth}{llX}
\toprule
\textbf{Endpoint} & \textbf{Method} & \textbf{Purpose} \\
\midrule
\texttt{/health}               & GET  & Health check \\
\texttt{/topics}               & GET  & Return seven predefined civic topic strings \\
\texttt{/avatar}               & GET  & Load canonical avatar from \texttt{avatars/canonical.json} \\
\texttt{/avatar/generate-image}& POST & Generate avatar image via Imagen \\
\texttt{/avatar/save}          & POST & Persist canonical avatar to disk \\
\texttt{/generate/image}       & POST & Single image generation \\
\texttt{/generate/post}        & POST & Full post pipeline (Agents 1, 2, Creative Storyteller) \\
\texttt{/generate/video\_async}& POST & Launch background video job (Agents 3, 4); returns job ID \\
\texttt{/jobs/\{job\_id\}}     & GET  & Poll job status (queued / running / done / error) \\
\bottomrule
\end{tabularx}
\end{center}

Video jobs run in background threads managed by a thread-safe job dictionary
(\texttt{\_jobs}, \texttt{\_jobs\_lock}). A 35-second timeout on post rendering
prevents the frontend from hanging. CORS is configured to allow \texttt{localhost:5173}
and \texttt{127.0.0.1:5173} by default, overridable via the \texttt{CORS\_ALLOW\_ORIGINS}
environment variable.

\subsection{Frontend: React Application}

The frontend is a React 19 + TypeScript 6.0 + Vite 8.0 single-page application.
It is organized into two tabs: \textbf{Generate} and \textbf{Avatar}.

The Generate tab exposes the full pipeline to the user. A topic selector populates
from the \texttt{/topics} endpoint; a platform toggle switches between LinkedIn and
Instagram; an audience personalization sidebar lets users configure segment, tone,
and platform style (stored in \texttt{localStorage}). Phase tracking (idle $\to$
post running $\to$ post done $\to$ video running $\to$ video done) drives the UI
state. A storyboard editor lets users review and edit each of the five scenes before
video generation is triggered. Once a video job is submitted, the frontend polls
\texttt{/jobs/\{job\_id\}} every 2--3 seconds and offers an MP4 download once the
job completes.

The Avatar tab provides a text description editor for a customizable AI persona, an
image generation preview, and a save action that persists the canonical avatar to the
backend. The avatar description and image are injected into Agent 2 and Creative
Storyteller prompts when present, ensuring visual consistency across posts.

All pipeline state (retrieval pack, content draft, final output, video output, QC
result) is persisted in \texttt{localStorage}, allowing non-volatile phase recovery
across page reloads without re-running the pipeline.

\subsection{Publishing Infrastructure}

A lightweight publishing queue (\texttt{publishing/service.py}) supports draft,
queued, and published states with optional scheduled publish times. A webhook
notification helper (\texttt{publishing/webhook.py}) POSTs a
\texttt{publish\_attempt} event payload to a configurable URL. LinkedIn and Instagram
publishing adapters (\texttt{publishing/adapters.py}) are currently stub
implementations; live platform API integration is planned but not yet active.

\subsection{Stack and Configuration}

\textbf{Backend dependencies} (selected): \texttt{google-generativeai $\geq$ 0.8.0},
\texttt{google-cloud-aiplatform $\geq$ 1.70.0}, \texttt{google-cloud-texttospeech
$\geq$ 2.16.0}, \texttt{moviepy $\geq$ 2.0.0}, \texttt{feedparser $\geq$ 6.0.11},
\texttt{trafilatura $\geq$ 1.12.0}, \texttt{fastapi}, \texttt{uvicorn},
\texttt{python-dotenv}.

\textbf{Frontend dependencies}: React 19.2.5, TypeScript 6.0.2, Vite 8.0.9.

\textbf{Required environment variables} (stored in \texttt{.env}, never committed):
\texttt{GOOGLE\_API\_KEY} (Gemini and embedding API access),
\texttt{GOOGLE\_CLOUD\_PROJECT} (default: \texttt{interstudent-nyc-2026}),
\texttt{GOOGLE\_CLOUD\_LOCATION} (default: \texttt{us-central1}),
\texttt{NYC\_OPEN\_DATA\_APP\_TOKEN} (optional, raises rate limits).

\textbf{To run locally:}
\begin{lstlisting}
# Backend
pip install -r requirements.txt
cp .env.example .env   # fill in API keys
python -m uvicorn backend.server:app --reload --port 8000

# Frontend
cd frontend && npm install
cp .env.example .env   # set VITE_API_BASE=http://localhost:8000
npm run dev
\end{lstlisting}

\textbf{Repository:} \url{https://github.com/gerald512-beep/InterStudent}

\newpage

% ─────────────────────────────────────────────────────────────────────────────
\section{Results and Analysis}
% ─────────────────────────────────────────────────────────────────────────────

\subsection{What the Pipeline Produces}

A full pipeline run on the topic ``Opening a US bank account as an international
student'' produces:

\begin{itemize}[leftmargin=1.5em]
  \item A 150--250-word LinkedIn or Instagram post with hashtags, a call-to-action,
    and cited source URLs drawn from live Google Search results and financial RSS feeds.
  \item A 1024$\times$1024 social image generated by Imagen depicting the avatar
    persona in a context relevant to the post topic.
  \item A five-scene 9:16 portrait video (approximately 40 seconds at 8 seconds per
    scene), with scene-specific visuals generated by Veo 3.1 Fast, a Chirp3-HD
    voiceover, and subtitle overlays.
  \item A quality control report scoring the content across seven criteria, with a
    binary Publish / Modify decision and per-criterion feedback.
\end{itemize}

Total wall time for a full run on tested hardware is approximately 4--6 minutes: roughly
15--20 seconds for Agents 1 and 2 combined, 30--60 seconds for the Creative Storyteller
image pass, and 3--4 minutes for parallel Veo video generation. The frontend handles
the post and video phases asynchronously, so users see post content within one minute
while the video processes in the background.

\subsection{What Works}

\textbf{Source retrieval quality.} Google Search Grounding consistently returns timely,
cited results on financial and civic topics that are more current than anything in the
static RSS or NYC Open Data feeds. This makes Agent 1's grounding-first design the
correct priority ordering: the live web is the most reliable signal, with the structured
datasets providing local context.

\textbf{Agent 2 content quality.} Posts generated for well-defined topics (bank accounts,
tax filing, student loan options) are generally readable, platform-appropriate, and
substantively accurate at the level of broad guidance. The persona configuration
(\texttt{config/persona.py}) successfully shapes tone across LinkedIn (professional,
informational) and Instagram (warmer, CTA-forward) without requiring separate prompts
per platform.

\textbf{Avatar persona consistency.} The canonical avatar system works as designed.
Injecting the avatar description into both Agent 2 and the Creative Storyteller produces
visually consistent persona representation across posts and storyboard scenes, which
is critical for building a recognizable channel identity.

\textbf{Video pipeline robustness.} The Veo $\to$ Imagen $\to$ Ken Burns fallback chain
handles API failures gracefully. In testing, fallback to Imagen with a Ken Burns zoom
effect produced acceptable video quality when Veo was unavailable, and the Chirp3-HD
audio track was indistinguishable between Veo and Imagen fallback paths because the
audio is replaced in all cases.

\textbf{Agent 4 decision quality.} The quality control agent correctly identified posts
that were too generic or that contained hedged language inappropriate for a public-facing
social channel, producing actionable per-criterion feedback rather than a binary pass/fail.

\subsection{What Does Not Work Well}

\textbf{Sensitive topic handling.} The current implementation has no explicit guardrails
for content on immigration compliance or housing rights beyond Agent 4's compliance
criterion. A post flagged as compliant by Agent 4 may still contain implicit guidance
that a student could misinterpret as legal advice. This is the highest-risk gap in the
current system.

\textbf{Single persona.} The single \texttt{PERSONA\_CONFIG} is calibrated for a broadly
``East Asian female student, 20--28 years old'' persona. This is a reasonable hackathon
simplification but limits the system's ability to serve the full diversity of the
international student population. Students from South Asia, Latin America, or Sub-Saharan
Africa face meaningfully different civic information challenges, and a single persona
configuration will not accurately represent or serve them all.

\textbf{In-memory vector store.} The absence of a persistent vector database means that
every run re-embeds all retrieved documents from scratch. For a hackathon demo this is
acceptable; for a production system handling dozens of topics across dozens of runs per
day, this is both slow and expensive.

\textbf{Publishing stubs.} LinkedIn and Instagram adapters are placeholders. The
publishing queue correctly manages draft and scheduled states, but no content can
actually be published to a live platform without activating the adapter layer.

\textbf{AI-sloppy behavior.} On broad or ambiguous topics (e.g., ``city safety for
students''), Agent 2 sometimes produces content that is structurally well-formed but
substantively vague---generic safety tips that any resident could find on any city
website, with no student-specific framing. This is the ``content quality theater'' risk
the slide deck identifies, and Agent 4 does not consistently catch it.

\subsection{Informal Testing}

The primary testing mode for this project was end-to-end pipeline execution: running the
full Agent 1 $\to$ Agent 2 $\to$ Creative Storyteller $\to$ Agent 3 $\to$ Agent 4 sequence
and evaluating the final video output for coherence, accuracy, and visual quality. This
form of integration testing is appropriate given the nature of the system---the most
meaningful unit of output is the complete content package, not any individual agent's
intermediate result.

The central challenge identified through this testing is pipeline stability and
consistency. Each run involves six distinct generative model calls across four different
model families (Gemini, Imagen, Veo, Chirp3-HD), each of which introduces variance.
A post that passes Agent 4's quality threshold on one run may produce a noticeably
different video on a second run with identical inputs, due to the stochastic nature of
Veo clip generation and Gemini's content drafting. Achieving consistent output quality
at scale---where a channel operator expects a predictable, on-brand result from every
pipeline run---remains the primary engineering challenge before this system is production-ready.

The team ran the pipeline across all seven predefined topics and spot-checked outputs
for factual accuracy against the cited source URLs. Posts on tax filing and banking
showed the highest factual alignment with sources. Posts on housing and safety showed
the highest rate of hedged or generic content.

\subsection{Economics}

The following cost estimates are approximate and based on public API pricing at time of
writing. Actual costs vary with token counts, image resolution, and video length.

\begin{center}
\begin{tabularx}{\textwidth}{Xll}
\toprule
\textbf{Component} & \textbf{Approx. Cost per Run} & \textbf{Notes} \\
\midrule
Gemini 2.5 Flash (Agents 1, 2, 4, CS) & \$0.01--\$0.05 & Depends on context length \\
\texttt{gemini-embedding-001}          & \$0.001--\$0.005 & Per retrieval pass \\
Google Search Grounding               & \$0.03--\$0.05 & Per grounded query \\
Imagen 4 Ultra (image gen)            & \$0.04 & Per image \\
Veo 3.1 Fast (5 scenes, $\sim$40s)    & \$0.50--\$0.75 & \$0.10--\$0.15/sec on Vertex AI \\
Chirp3-HD TTS                         & \$0.01--\$0.02 & Per minute of audio \\
NYC Open Data                         & \$0.00 & Free API \\
\midrule
\textbf{Estimated total per full run} & \textbf{\$0.10--\$0.20} & Excl. Veo \\
\bottomrule
\end{tabularx}
\end{center}

At \$0.10--\$0.20 per full content run (excluding Veo, which is in enterprise preview),
a student organization producing five posts per week would spend approximately
\$3--\$5 per month in API costs. This is well within the budget of a university
department or student organization. If offered as a SaaS product, a subscription in
the \$30--\$50 per month range for small organizations would provide a healthy margin
while remaining accessible.

\subsection{Competitive Landscape}

The project occupies a niche that existing tools do not serve well.

\begin{itemize}[leftmargin=1.5em]
  \item \textbf{General-purpose AI assistants} (ChatGPT, Gemini, Claude): Can answer
    individual civic questions but do not monitor live data, produce formatted posts,
    or maintain a consistent persona and publishing schedule. No civic data integration.
  \item \textbf{Social media management platforms} (Buffer, Hootsuite, Sprout Social):
    Excellent scheduling, analytics, and team workflow tools, but no native ability to
    generate civic content from raw data sources. Content must be written by a human
    and then scheduled.
  \item \textbf{AI content generators} (Jasper, Copy.ai, Writesonic): Generate
    marketing copy effectively but are not designed for civic information, do not
    integrate structured government data sources, and have no multimodal video pipeline.
  \item \textbf{University international student office portals}: Authoritative and
    trustworthy, but static, text-heavy, not distributed through social platforms, and
    not capable of responding to live policy or data changes in near real time.
\end{itemize}

The project's differentiator is the combination of live, structured civic data
integration, domain-specific persona design, and full multimodal output (text + image
+ video) from a single pipeline run. No comparable tool targeting this specific
audience and use case was identified in the competitive survey.

\newpage

% ─────────────────────────────────────────────────────────────────────────────
\section{Conclusion}
% ─────────────────────────────────────────────────────────────────────────────

\subsection{Main Findings and Lessons}

The International Student AI Influencer demonstrates that a coordinated set of
specialized AI agents can bridge the gap between raw civic data and polished,
student-appropriate social media content at low marginal cost. The key architectural
lesson is that specialization pays: giving each agent a narrowly defined responsibility
(retrieve, write, render, produce video, evaluate) produces better output than asking
a single model to do everything in one prompt. The pipeline is also more debuggable
as a result---failures and quality problems can be traced to a specific agent rather
than to an opaque monolithic prompt.

The most important honest finding is that technical quality and informational trustworthiness
are not the same thing. The pipeline can produce content that is visually polished,
grammatically correct, and structurally on-brand while being subtly wrong or too
generic to be actionable. This gap is not solvable by making the models larger or the
prompts more sophisticated---it requires the human approval gate and user testing called
for in the roadmap.

\subsection{Extensions and Roadmap}

The immediate next steps, as defined in the project roadmap, are:

\begin{itemize}[leftmargin=1.5em]
  \item Source citations are already embedded in generated posts and displayed to the
    user in the UI; the next step is ensuring citations appear consistently across all
    output formats, including the video storyboard and published captions.
  \item Implement a human approval gate in the publishing queue, requiring explicit
    sign-off before any content moves from draft to live.
  \item Conduct structured user testing with real international students, measuring
    both content accuracy (fact-checking against sources) and perceived relevance
    (student survey).
  \item Build explicit content filters for sensitive domains (housing rights,
    immigration compliance, safety guidance) that trigger additional human review or
    disclaimer insertion.
  \item Activate LinkedIn and Instagram publishing adapters to close the gap between
    queue management and live distribution.
  \item Expand platform support to TikTok, WhatsApp, and email newsletters, each of
    which requires format and tone adaptation beyond the current LinkedIn/Instagram toggle.
\end{itemize}

Looking 12--24 months ahead, two developments in the AI landscape are most likely to
change what this product can do. First, continued improvements in video generation
models (Veo and its successors) will reduce per-scene generation time and improve lip
sync accuracy, making the video pipeline practical to run at higher volumes. Second,
improved grounding and citation frameworks will make it easier to build a verifiable
chain from a published claim back to its authoritative source, which is the core
requirement for trust in civic information content. A production version of this system
should be built on that foundation---where every claim in a post carries a checkable
source link, not merely a cited URL from a retrieval result.

\subsection{Closing}

The project establishes a technically credible foundation for a trust-ready civic AI
influencer pipeline. It surfaces the right problems, builds an architecture that can
accommodate the solutions, and is honest about the gap between what the prototype
can do today and what a responsible production deployment would require. The path
forward is clear: source verification, human oversight, and real student testing. Those
three steps, more than any further model capability, are what would make this tool
genuinely useful to the students it is designed to serve.

\newpage

% ─────────────────────────────────────────────────────────────────────────────
\section*{References}
\addcontentsline{toc}{section}{References}
% ─────────────────────────────────────────────────────────────────────────────

\begin{enumerate}[leftmargin=2em, label={[\arabic*]}]

  \item Google DeepMind. (2025). \textit{Gemini 2.5 Flash model documentation}.
    Google AI for Developers.
    \url{https://ai.google.dev/gemini-api/docs}

  \item Google DeepMind. (2025). \textit{Imagen 4 on Vertex AI}.
    Google Cloud Documentation.
    \url{https://cloud.google.com/vertex-ai/generative-ai/docs/image/overview}

  \item Google DeepMind. (2025). \textit{Veo on Vertex AI}.
    Google Cloud Documentation.
    \url{https://cloud.google.com/vertex-ai/generative-ai/docs/video/overview}

  \item Google Cloud. (2025). \textit{Text-to-Speech: Chirp3-HD}.
    Google Cloud Documentation.
    \url{https://cloud.google.com/text-to-speech/docs}

  \item NYC Open Data. (2025). \textit{NYC Jobs} (Dataset \texttt{kpav-sd4t}).
    City of New York.
    \url{https://data.cityofnewyork.us}

  \item NYC Open Data. (2025). \textit{CUNY Enrollment} (Dataset \texttt{d44m-4xgv}).
    City of New York.
    \url{https://data.cityofnewyork.us}

  \item Internal Revenue Service. (2025). \textit{Taxation of Nonresident Aliens (Publication 519)}.
    \url{https://www.irs.gov/publications/p519}

  \item U.S. Citizenship and Immigration Services. (2025). \textit{Students and Exchange
    Visitors: F-1 and J-1 Visa Overview}.
    \url{https://www.uscis.gov}

  \item NASFAA. (2025). \textit{Financial Aid for International Students}.
    National Association of Student Financial Aid Administrators.
    \url{https://www.nasfaa.org}

  \item FastAPI. (2025). \textit{FastAPI documentation}.
    \url{https://fastapi.tiangolo.com}

  \item Ziadé, T., \& Stinner, C. (2025). \textit{MoviePy documentation}.
    \url{https://zulko.github.io/moviepy}

  \item OpenAI. (2023). \textit{ChatGPT: Optimizing Language Models for Dialogue}.
    OpenAI Blog. \emph{[Referenced for competitive analysis.]}

  \item Buffer Inc. (2025). \textit{Buffer: Social Media Management}.
    \url{https://buffer.com}

  \item Hootsuite Inc. (2025). \textit{Hootsuite: Social Media Management}.
    \url{https://hootsuite.com}

\end{enumerate}

\newpage

% ─────────────────────────────────────────────────────────────────────────────
\appendix
\section{Environment Variable Reference}
% ─────────────────────────────────────────────────────────────────────────────

\begin{center}
\begin{tabularx}{\textwidth}{lXl}
\toprule
\textbf{Variable} & \textbf{Purpose} & \textbf{Required} \\
\midrule
\texttt{GOOGLE\_API\_KEY}          & Gemini and embedding API access            & Yes \\
\texttt{GOOGLE\_CLOUD\_PROJECT}    & Vertex AI project ID                       & Yes \\
\texttt{GOOGLE\_CLOUD\_LOCATION}   & Vertex AI region (default: us-central1)    & Yes \\
\texttt{NYC\_OPEN\_DATA\_APP\_TOKEN} & Socrata API token (raises rate limits)   & No \\
\texttt{CORS\_ALLOW\_ORIGINS}      & Override default CORS whitelist            & No \\
\texttt{VITE\_API\_BASE}           & Frontend API base URL                      & Yes (frontend) \\
\bottomrule
\end{tabularx}
\end{center}

% ─────────────────────────────────────────────────────────────────────────────
\section{Predefined Topic List}
% ─────────────────────────────────────────────────────────────────────────────

The backend \texttt{/topics} endpoint returns the following seven topics at startup:

\begin{enumerate}[leftmargin=2em]
  \item Opening a US bank account as an international student
  \item Understanding your F-1 visa employment restrictions
  \item Filing US taxes as a non-resident alien
  \item Finding affordable housing in NYC as a student
  \item Navigating NYC public transportation and city services
  \item Student loan options for international students
  \item Understanding your rights as an international student in the US
\end{enumerate}

% ─────────────────────────────────────────────────────────────────────────────
\section{Agent Quality Control Scoring Rubric}
% ─────────────────────────────────────────────────────────────────────────────

Agent 4 scores output against the following seven criteria, each on a 1--10 scale.
An aggregate score at or above 7 triggers a \texttt{PUBLISH} decision.

\begin{center}
\begin{tabularx}{\textwidth}{lX}
\toprule
\textbf{Criterion} & \textbf{What is evaluated} \\
\midrule
Accuracy       & Factual alignment between post claims and cited sources \\
Tone           & Appropriateness of register and voice for the target audience \\
Engagement     & Predicted reader interest, CTA clarity, and hashtag relevance \\
Compliance     & Absence of implicit legal, financial, or immigration advice \\
Platform fit   & Formatting, length, and structure appropriate to LinkedIn or Instagram \\
Video script   & Coherence and flow of the SSML voice script across five scenes \\
Video visual   & Alignment between visual prompts and the post narrative \\
\bottomrule
\end{tabularx}
\end{center}

\end{document}
